% ============================================================
% Chapter 6: Experimental Design
%   Section: Selection and Justification of the Evaluation Dataset
%     Subsection: Why FEVER Is Not Suitable for This Work
% Included from chapters/06-experimental-design/02-evaluation-dataset/evaluation-dataset.tex
% ============================================================

\subsection{Why FEVER Is Not Suitable for This Work}
\label{subsec:discard-justification}

The system developed in this thesis --- described in detail in
Chapter~\ref{ch:methodology} --- does not perform verification
against a closed corpus. Its architecture is organised into six
sequential stages:

\begin{enumerate}
    \item \textbf{Admission}: ingestion and decomposition of the
    news article into atomic, verifiable claims.
    \item \textbf{Selection}: entity extraction and formulation of
    search queries from each claim.
    \item \textbf{Retrieval}: retrieval of candidate documents via
    SearXNG, a metasearch engine that queries multiple search
    engines over the open web in real time.
    \item \textbf{Ranking}: ordering of the retrieved sources by
    relevance to the claim.
    \item \textbf{LLM check}: reasoning over the ranked evidence by a
    language model (Llama 3.1, run locally via Ollama) to produce a
    preliminary verdict.
    \item \textbf{Recalibration}: adjustment of the verdict and its
    associated confidence, propagating the result to the article's
    overall verdict.
\end{enumerate}

This architecture --- referred to throughout this work as the
\textbf{proposed method} --- therefore performs
\textbf{open-domain retrieval}: no fixed corpus exists against which
claims are compared; instead, evidence is searched for \textit{on
the fly} on the live web at verification time, exactly as a human
fact-checker would.

Evaluating this system against FEVER or its multilingual variants
would introduce a serious methodological mismatch, for two reasons:

\begin{itemize}
    \item \textbf{Invalid-evidence bias.} The system could locate,
    through SearXNG, evidence that is perfectly valid and sufficient
    in a source other than the predefined Wikipedia sentence, and
    still be penalised for not matching the expected evidence
    verbatim. The resulting metric would systematically underestimate
    the system's true performance.
    \item \textbf{Partial evaluation of the architecture.} Since
    FEVER fixes the source corpus in advance, evaluation against this
    benchmark would measure almost exclusively the \textit{LLM check}
    stage, leaving the \textit{Selection}, \textit{Retrieval}, and
    \textit{Ranking} stages unevaluated --- precisely the stages that
    constitute the most original, and also the most fragile,
    contribution of the system, as discussed in
    Chapter~\ref{ch:results}.
\end{itemize}

For these reasons, FEVER/XFEVER was discarded as the primary
evaluation benchmark: not out of convenience or expected performance,
but because the dataset's own design assumes a (closed) retrieval
paradigm incompatible with the (open) retrieval paradigm underlying
the proposed system.
